\chapter[Sermon 64. Another Angel Preaches, That Babylon Shall Fall: and Another Dissuades All Men from the Fellowship of the Religion of the Beast.]{Another Angel Preaches, That Babylon Shall Fall: and Another Dissuades All Men from the Fellowship of the Religion of the Beast.}
\label{sermon:064}
\markright{Sermon 64. Another Angel Preaches, That Babylon Shall Fall: and ...}

And there followed another angel, saying: “She has fallen, she has fallen—Babylon the great city—because she has made all nations drink of the wine of her fornication.” And the third angel followed them, saying with a loud voice: “If anyone worships the beast and his image, and receives his mark on the forehead or on the hand, that person shall drink of the wine of the wrath of God, which is poured out in the cup of his wrath. And that person shall be punished with fire and brimstone before the holy angels and before the Lamb, and the smoke of their torment ascends forever. They have no rest day or night, those who worship the beast and his image, and whoever receives the imprint of his name.” Here is the perseverance of the saints. Here are those who keep the commandments and the faith of Jesus.

\index{Antichrist}For the comfort of the faithful flock of Christ, another angel appears, a type of all godly preachers, who proclaims with great constancy that the kingdom of Antichrist will fall, however it promises itself everlastingness. From this, the saints gather that persecutions will certainly end, along with all other abominations throughout the world. For, considering the continual persecutions of the wicked, the saints cannot but be marvelously saddened; they must needs receive no small joy and comfort from this understanding, that they here will not endure forever.

\index{Apocalypse (Book of)}\index{John, Saint}\index{Jerome, Saint}\index{Tertullian}\index{Rome}and here it is said that Babylon shall fall. And in deed it were foolishness to expound these things of the old Babylon in Asia, which was fallen long since, scarcely any token thereof being left; we must therefore understand it of another, which is in her flower, and even by a figurative speech, we must understand it of Rome. For there is a great correspondence as it were, betwixt both: Babylon was the first monarchy, Rome is the last. Babylon sore afflicted the people of God, so doth Rome grievously vex the church of God. Babylon burdened Israel with a grievous captivity: so Rome vexes the church, with more than a long captivity. Babylon overcame the people of God, and burning the City of Jerusalem, and destroying the temple, led away Israel captive: so Rome also having razed the city of Jerusalem, and subverting the temple, triumphed of Israel. Babylon planted idolatry, superstition, and all abomination, advanced, maintained, and set forth the same unto all men: but at the length when she would have thought least of it, the people of God being suddenly delivered, she was utterly subverted. So is Rome also, the mother and nurse, and reviver of all abominations in the church of the last time, wherein she shall perish at the last, all those that believe truly in Christ being delivered. And especially it is called great. For how great and mighty the church of Rome is, all we see and by experience know at this day. Neither am I the first that understand by Babylon, Rome. For many expositors reading the first Epistle of Saint Peter, in the end of the epistle, do understand by Babylon, Rome. Certainly Oecumenius says: And here he calls Babylon Rome, for the excellency and brightness of the Empire: the which Rome obtained a long time since. But this the more ancient writers expound more plainly, as Tertullian in his book against the Jews, which says: so Babylon with Saint John bears the figure of the City of Rome, therefore also great, and proud in her kingdom, and a murderer of the saints. The same words in a manner, he repeats in the third book against Marcion. And no less plainly Saint Jerome calls Rome Babylon: and that same Babylon whereof S. John speaks in the Apocalypse. Read the epistle of Paula and Eustochium written to Marcella, by the help of Saint Jerome. Read himself in the 11th question to Algasia. Again in the preface to the book of Didymus of the Holy Ghost, to Pauliniane. Also in the end of the 2nd book against Jovinian. The same in the life of Saint Mark: Peter, saith he, in the first epistle, under the name of Babylon, doth figuratively signify Rome. But Saint John will expound himself in the 17th chapter. And we understand that the City of Rome shall fall chiefly, with all her ungodliness: and with the same also, the Roman superstition and abomination, throughout the world. And the Angel in deed, says she is fallen, which is yet to fall: and that by the prophetical manner of speaking, wherein that which shall assuredly come to pass is uttered, as though it were now done. To signify the certainty thereof, that reduplication or iterating of the word also appertains, she is fallen, she is fallen. This is also repeated in the 18th chapter, where it shall be showed how it is taken out of the Prophets, \&c. Notwithstanding both a desire and joy also, might seem here to be signified. For such things as we have long, and with a desire looked for, we receive now coming and say, thou art come, thou art come at the last long looked for, and now make me glad. For the saints with a great desire, look and long for the destruction of that most wicked, most vile, and most troublesome kingdom of Antichrist.

\index{Justification}The cause of the destruction of the commonwealth and the church of Rome is also shown, for she has caused all nations to drink, and has made them drunken with the wine of wrath of her fornication. Indeed, the effect of wine on people is greatest. Therefore, doctrine is compared to it in the Prophets. Thus, Rome, with her unclean and corrupt opinions, has made all people drunken. And it is called the wine of wrath. For, as God is angry with those who err in the way of the Roman church, he allows them to do so. Because God has revealed the sincere doctrine of life through his only Son and most chosen Apostles, and people do not receive it, God is justly offended with them and gives them over to a reprobate mind, that they may follow shameful errors, as Saint Paul prophesied would come to pass in 2 Thessalonians, chapter 2. This wine is also called the wine of her fornication, by which she herself, having first played the harlot, has now become the master of fornication, and as it were a bawd to all others. This manner of speaking is well known even from the Prophets. Rome did not persist in the doctrine of the Gospel and of the Apostles, but invented a new one, and that contrary to the Gospel, concerning the vicar of Christ on earth, of the power of keys, of indulgences and pardons, of justification by works and merits, of satisfaction and confessions, of the worship of images, and praying to saints, of celebrating masses, and worshipping the sacrament of the altar, as they term it, of monasticism and vows, and such other innumerable things. This doctrine, as she presents it as Apostolic, ancient, and Christian, she offers to all people, and so plucks them from Christ, withdraws them from the Gospel, seduces them from the old Christianity, and destroys innumerable souls. Therefore, God pours out to her also from the cup of his wrath, and brings her also to destruction forever.

\index{Papacy}And on this occasion he dissuades all men from fellowship with the Roman church or papistry, that we have nothing to do with the Roman religion unless we will be partakers also of the everlasting punishment. He reasons therefore of the loss and punishments, and describes grievous and horrible pains, if perhaps men might so be frightened from that ungodliness. The angel therefore cries, and that with a loud voice. Wherefore let all ecclesiastical preachers learn that they must earnestly and terribly cry out in this case, that all flee the communion of the Roman or popish church. I know doubtless what the common people believe and say, that all shall be saved at the last day, whatever religion they be of; and especially if any remain an open papist. But we can neither condemn nor absolve any man, set them in Heaven, or cast them to Hell. God is a righteous judge. He alone knows who shall be saved or damned. We ought therefore of right to credit his judgments. But whereas he pronounces openly that the favorers of the Roman church shall be damned, who am I to say the contrary, or what men will pronounce otherwise? Let us hear therefore the sentence of the just judge, and let us believe the word of the Son of God, and let us beware of the popish religion.

\index{Arethas of Caesarea}\index{Repentance}I have already explained in chapter 13 what it means to worship the beast and his image, and what it means to receive the mark on the forehead and on the right hand. Briefly, they worship and receive the mark of the beast, those who participate with the Roman church or religion; finally, those who obey the wicked decrees of the empire and persevere in obedience to the See without repentance. Arethas, interpreting this passage, says that to worship the beast and to receive his seal is to consider Antichrist to be God, and in word and deed to promote the things he desires.

And here in a horrifying manner, and with prophetic words, is described everlasting damnation, prepared for those who, forsaking Christ the Savior, cleave unto Antichrist the destroyer. Like as they have drunk of the corrupt doctrine infused by the Pope, so again shall they drink, that the just Lord shall pour out of the cup of wrath. And the wine that is poured into the cup of God’s wrath is the straight, exquisite, and most grievous judgment of God, wherein, being angry, he inflicts upon the Antichristians horrible and unspeakable punishment. A like manner of speech is read in Jeremiah, chapter 25. And like as pure wine, not delayed, is of most efficacy and pierces: so the judgment of God, wherein he will proceed against the Antichristians, shall be most grievous, such as no tongue, be it never so eloquent, can express.

And for a further declaration shortly after follows what they must drink: truly fire and brimstone. Perhaps the Lord alluded to these words of David in the eleventh Psalm: Upon the ungodly he shall rain snares, fire, brimstone, storm, and tempest; this reward they shall have to drink. He seems moreover to have alluded to the burning of Sodom, and to the thirtieth chapter of Isaiah, in the end of which it is shown that hell shall be wide enough to receive all the ungodly, and that matter shall never want to nourish the fire, nor the fire be quenched. He expresses moreover a grievous pain, where he says that they shall be tormented, and that in the sight of the Lamb and holy angels, that so they may receive deserved punishment forever for their contempt, whereby they have despised the Lamb and messages of angels. Likewise in Luke 13, the Lord says: There shall be weeping and gnashing of teeth, when they shall see Abraham, Isaac, and Jacob, and all the prophets in the kingdom of God, and you to be shut out, and so forth.

And that same pertains also unto evidence, and to stir up a terror in the minds of all men, where he adds by a figurative speech: and the smoke of their torment ascends up evermore. Therefore, the burning and punishment of the ungodly will be everlasting, and never to be finished world without end. And we seem here at this description, as it were before our eyes to see the flames of eternal damnation carried up on high, and cast up with them great heaps of smoke, to roll up and disperse them far and wide. I remember here that of Virgil.

And that no kind of terror might be lacking, he expresses most aptly and most abundantly the perpetuity of everlasting punishment, saying: neither have they rest day nor night. So says the Lord in Mark 9. Their fire is never quenched, and their worm shall never die. They err, therefore, who promise to the damned deliverance from their torments after many worlds.

\index{Hebrews, Epistle to the}And it is not in vain that he repeats what he had said before, how those who worship the beast shall suffer these things. And therefore he repeats it, lest, as it happened, we should consider it a trivial matter. They shall be condemned, says the truth, who receive the Papist cult and religion, and persevere in the same. To all this is annexed an acclamation, or double sentence, notable and wholesome. For inasmuch as the wisdom of God did foresee what adversity remained for the godly in this world, which they might surely expect at Antichrist’s hand, who professed the truth, therefore for a comfort and consolation he adds: here is the patience of the saints, which is as much as if he had said: and here shall patience take place, whereby the saints may overcome all evils. Here we have need of stout courage, and a sure and constant mind. In Luke 12, the Lord likewise requires patience in persecutions. Here therefore is counsel given, how the saints should behave themselves, namely that they should suffer patiently those evils that Antichrist shall work against them. And there follows another sentence, which clarifies this: here are they that keep the commandments of God, and the faith of Jesus. They shall overcome through patience in so great evils and dangers, which keep the commandments of God, the foundation of which is the faith of Jesus Christ; which truly put all their trust in Christ, hear the word of the gospel, and keep the commandments of God, not of men. The like unto these are read in Matthew 24 and in Hebrews 10. Arethas says in this same time of Antichrist, the patience of the saints is shown. Then the speech is figured, as it were by a question moved. And who are they whom he calls patient? After, as though he should answer: they that keep the commandments of God, and the faith of Jesus. For they when perils approach will value God more than death and temporal evils. This he says. I pray God these things be as faithfully performed by us, as they are easily understood. The Lord grant us his spirit.
